\documentclass{article}
\usepackage{geometry}
\geometry{tmargin=1in,bmargin=1in,lmargin=1in,rmargin=1in}
\usepackage{array}
\usepackage{float}
\usepackage{booktabs}
\usepackage{amsmath}
\allowdisplaybreaks
\usepackage{amssymb}
\usepackage{graphicx}
\usepackage{setspace}
\usepackage[normalem]{ulem}
\usepackage{graphicx}
\usepackage{caption}
\usepackage[flushleft]{threeparttable}
\usepackage{longtable}
\usepackage{multirow}
\usepackage{float}
\usepackage{mathtools}
\usepackage{cancel}
\usepackage{comment}
\raggedbottom
\onehalfspacing
\usepackage{lscape}
\usepackage{xcolor}
\usepackage{comment}
\usepackage{pdflscape} 
\usepackage{ulem}
\PassOptionsToPackage{hyphens}{url}\usepackage{hyperref}


\hypersetup{colorlinks, citecolor=blue,urlcolor=blue,allcolors=blue}
\newcommand\numberthis{\addtocounter{equation}{1}\tag{\theequation}}


\begin{document}
	\noindent
	Notes on incorporating new AviationDB data into pilot study\\
	Xinyu \\
	03-31-26

\section{Motivational Evidence}
\subsection{Pilots are Mobile}
\begin{itemize}
	\item Airline transport pilots and commerical pilots have much higher interstate migration rates (4-7\%) compared to the general population (2-3\%)
\end{itemize}
	
	\begin{figure} [H]
		\centering
		\includegraphics[scale=0.6]{../output/aviationdb/figures/migration_rates_by_pilot_group.png}
	\end{figure}
	
	
\section{Correlating Taxes with Pilot Density}

\begin{itemize}
	\item Binscatter plots showing that lower tax states have a higher concentration of pilots, mainly driven by zero-income tax states
	\item x-axis: average state tax rate (top), average state + federal tax rate (bottom)
	\item y-axis: Share of pilots with respect to total working population in each state by year
\end{itemize}
\begin{figure} [H]
	\centering
	\includegraphics[scale=0.45]{../output/aviationdb/figures/binscatter_astr_pilots_bls_airline_mean_no_AK_HI.png}
\end{figure}

\begin{figure} [H]
	\centering
	\includegraphics[scale=0.45]{../output/aviationdb/figures/binscatter_atr_pilots_bls_airline_mean_no_AK_HI.png}
\end{figure}


\section{Empirical Strategy}
\subsection{Stock Model (Agrawal et. al 2025)}
\textbf{Goal}: Do states with lower taxes have a larger population of pilots?

	\begin{equation*}
	\ln \biggr ( \frac{N_{s, t}}{N_{s', t}} \biggr)=  \beta \ln \biggr(\frac{1 - \tau_{s, t}}{1 - \tau_{s', t}} \biggr) + \gamma_s + \gamma_t +  u_{s, t}
\end{equation*}
where
\medskip
\begin{itemize}
	\small
	\item $s$ indexes state, $s'$ indexes the reference state, and $t$ indexes year.
	\medskip
	\item LHS:  population of pilots in state $s$ relative to the reference state
	\medskip
	\item RHS: net-of-income-tax rate in state $s$ relative to the reference state
	\medskip
	\item $ \gamma_s$ and $\gamma_t$:  state and year fixed effects
\end{itemize} 


\begin{table}[htbp]
   \caption{AFM stock model results: Balanced panel}
   \bigskip
   \centering
   \small
   \begin{tabular}{lccccc}
      \toprule
       & \multicolumn{5}{c}{log\_stock\_n}\\
                                 & BLS airline mean & BLS airline median & SOI p90       & SOI p95               & SOI p99 \\   
                                 & (1)              & (2)                & (3)           & (4)                   & (5)\\  
      \midrule 
      lnet\_atr                  & 2.518$^{**}$     & 2.414$^{*}$        & -0.5217       & -0.0466               & 1.099\\   
                                 & (1.184)          & (1.243)            & (1.607)       & (1.661)               & (0.7028)\\   
       \\
      Observations               & 441              & 441                & 441           & 441                   & 441\\  
      R$^2$                      & 0.99782          & 0.99782            & 0.99777       & 0.99777               & 0.99780\\  
      Within R$^2$               & 0.02319          & 0.02144            & 0.00111       & $6.88\times 10^{-6}$  & 0.01186\\  
       \\
      year fixed effects         & $\checkmark$     & $\checkmark$       & $\checkmark$  & $\checkmark$          & $\checkmark$\\   
      dest\_state fixed effects  & $\checkmark$     & $\checkmark$       & $\checkmark$  & $\checkmark$          & $\checkmark$\\   
      \bottomrule
   \end{tabular}
\end{table}




Notes:
\begin{itemize}
	\item The effects go away if I use the AGI thresholds from IRS SOI individual tax stats to construct the tax measures. 
\end{itemize}

\newpage
\begin{landscape}
	\subsection{Migration Log Odds Ratio (Moretti and Wilson 2017)}
	\textbf{Goal}: Do destination-origin state tax differentials affect the probability of moving from the destination state to the origin state? \\
	
	Regression:
	\begin{align*} 
		\underbrace{\log \Big(\frac{P_{odt}}{P_{oot}}\Big)}_{\substack{\text{migration}\\\text{log odds ratio}}}
		= \beta_0 + \beta_1 \underbrace{[\log(1 - \tau_{dt}) - \log(1 - \tau_{ot})]}_{\substack{\text{log dest-origin}\\\text{net-of-tax diff.}}}
		+ \underbrace{\gamma_t}_{\substack{\text{year}\\\text{FE}}}
		+ \underbrace{\gamma_d}_{\substack{\text{destination}\\\text{state FE}}}
		+ \underbrace{\gamma_o}_{\substack{\text{origin}\\\text{state FE}}}
		+ \underbrace{\gamma_{od}}_{\substack{\text{state-pair}\\\text{FE}}}
		+ u_{odt}
	\end{align*}
	where $o$ indexes the origin state, $d$ indexes the destination state, and $t$ indexes year.
	\begin{itemize}
		\item LHS: migration log odds ratio, between the number of pilots who moved from the origin state $o$ to the destination state $d$ and the number of pilots who stayed in the origin state $o$ in year $t$.
		\item RHS: log net-of-average-personal-income tax rate difference between the destination state and the origin state.
	\end{itemize}  
	
	\begin{figure} [H]
		\centering
		\includegraphics[scale=0.45]{../output/aviationdb/figures/mw_reg_binscatter_bls_airline_mean.png}
		\caption{}
	\end{figure}
	
	
\end{landscape}


\begin{table}[htbp]
   \caption{Moretti and Wilson (2017) Regression Specification Results, 3-way clustering: Balanced panel}
   \bigskip
   \centering
   \small
   \begin{tabular}{lccccc}
      \toprule
       & \multicolumn{5}{c}{lodds\_ratio}\\
                                                & BLS airline mean & BLS airline median & SOI p90              & SOI p95       & SOI p99 \\   
                                                & (1)              & (2)                & (3)                  & (4)           & (5)\\  
      \midrule 
      lnet\_tax\_diff                           & 6.744$^{***}$    & 7.296$^{***}$      & 0.6974               & 2.792         & 2.229$^{*}$\\   
                                                & (2.112)          & (2.137)            & (2.370)              & (2.584)       & (1.234)\\   
       \\
      Observations                              & 6,259            & 6,259              & 6,259                & 6,259         & 6,259\\  
      R$^2$                                     & 0.84948          & 0.84957            & 0.84885              & 0.84894       & 0.84903\\  
      Within R$^2$                              & 0.00422          & 0.00484            & $4.9\times 10^{-5}$  & 0.00064       & 0.00123\\  
       \\
      year fixed effects                        & $\checkmark$     & $\checkmark$       & $\checkmark$         & $\checkmark$  & $\checkmark$\\   
      origin\_state fixed effects               & $\checkmark$     & $\checkmark$       & $\checkmark$         & $\checkmark$  & $\checkmark$\\   
      dest\_state fixed effects                 & $\checkmark$     & $\checkmark$       & $\checkmark$         & $\checkmark$  & $\checkmark$\\   
      origin\_state-dest\_state fixed effects   & $\checkmark$     & $\checkmark$       & $\checkmark$         & $\checkmark$  & $\checkmark$\\   
      \bottomrule
   \end{tabular}
\end{table}





\subsection{Individual-Level Linear Probability Model}

\textbf{Goal:} Are pilots less likely to move out of the origin state if the net-of-tax rate is high? \\

Regression:
\begin{equation*}
	\text{Moved}_{i,o,t} =  \beta (1 - \tau_{o, t}) + \gamma_o + \gamma_t +  \gamma_i  +u_{i, o, t}
\end{equation*}
\begin{itemize}
	\item LHS: 1 if moved, 0 if not.
	\medskip
	\item RHS: net-of-ATR in origin state 
	\medskip
	\item $ \gamma_o$, $\gamma_t$, and $\gamma_i$:  origin state, year, and pilot fixed effects
\end{itemize} 



\begingroup
\centering
\small
\begin{tabular}{lcc}
   \tabularnewline \midrule \midrule
   Dependent Variable: & \multicolumn{2}{c}{moved}\\
   Model:             & (1)             & (2)\\  
   \midrule
   \emph{Variables}\\
   net\_origin\_atr   & -0.0098$^{***}$ & -0.0114$^{*}$\\   
                      & (0.0027)        & (0.0065)\\   
   \midrule
   \emph{Fixed-effects}\\
   year               & Yes             & Yes\\  
   origin\_state      & Yes             & Yes\\  
   unique\_id         &                 & Yes\\  
   \midrule
   \emph{Fit statistics}\\
   Observations       & 428,600         & 428,600\\  
   R$^2$              & 0.01395         & 0.20923\\  
   Within R$^2$       & 0.00013         & 0.00021\\  
   \midrule \midrule
   \multicolumn{3}{l}{\emph{Clustered (origin\_state \& unique\_id) standard-errors in parentheses}}\\
   \multicolumn{3}{l}{\emph{Signif. Codes: ***: 0.01, **: 0.05, *: 0.1}}\\
\end{tabular}
\par\endgroup




\subsection{Quasi-Experimental Design: Event Study}

This is still in progress. The idea is to do a DID utilizing the 2017 TCJA, or to do a staggered DID utilizing major state income tax changes (defined using a threshold, say 1pp). If successful, this should be the main empirical strategy.\\

\noindent
Current challenges:
\begin{itemize}
	\item There are only 5 major state income tax changes (above 1pp) during 2007-2022.
	\item Still need to figure out a reasonable control group for airline pilots. Previously the idea is to use commercial pilots as a control group, but based on the migration rates graph, commercial pilots seem to be similarly mobile, despite earning much less than airline pilots.
\end{itemize}
\end{document}










